\section{\textcolor{titleblue}{Introduction}}

RNA plays two main roles in cells. It can serve as an information
carrier, carrying the genetic code needed to make proteins, but it can also
serve a functional role such as catalyzing chemical reactions or regulating genes~\citep{doudna2002,breaker2012}.
Recent advances have shown that RNA has huge potential for developing new medicines.
For example, \citet{rojas2023} showed that personalized mRNA vaccines stimulated immune responses against tumor
targets in patients with pancreatic cancer. This advancement could open the path toward
cancer treatments that are tailored to individual patients. Additionally, on the computational side,
recent methods such as LinearDesign have shown that algorithmic techniques can be used
to design mRNAs with improved stability, protein production, and immune response~\citep{zhang2023}.
These advances in RNA technology motivate further work on designing RNA.

RNA design, also called inverse folding, is the task of finding a nucleotide
sequence that folds into a given target structure~\citep{hofacker1994,dirks2004}. Given a target structure,
a good RNA design method should find a sequence that favors the target
over the competing folds in its ensemble. However, this problem is difficult because
the number of possible sequences for a target grows exponentially with length,
and each of these sequences can adopt many alternative structures,
which makes the search especially hard for long RNAs~\citep{tang2026}. Existing methods that optimize
target probability face limitations in quality, runtime, or both. For example,
SAMFEO maintains several promising sequences and uses structure-aware mutations to optimize folding probability~\citep{zhou2023},
but it struggles on difficult and long targets, where the enormous design space
makes local search prone to getting stuck in local optima.
On the other hand, SamplingDesign~\citep{tang2026} produces high-quality designs, even for difficult puzzles in
Eterna100, but it is computationally expensive and too slow to be practical for
longer RNAs. Finding a high-probability sequence within a practical runtime remains a major challenge.

To address this bottleneck, we present EliteDesign, which uses the cross-entropy method~\citep{rubinstein1999,deboer2005} to
find RNA sequences with high target probability
in a computationally efficient way.
At each iteration, EliteDesign samples from the valid design space, scores the sequences,
and then picks a set of elite sequences based on target probability.
It then moves the current distribution toward the distribution described by the elite
set. During the sampling rounds, sequences are scored with LinearPartition~\citep{zhang2020},
which estimates folding probability in linear time, and the best final candidates are
rescored exactly with ViennaRNA~\citep{lorenz2011}.

We benchmark EliteDesign on several datasets and show that it finds higher target
folding probabilities than the methods we compare against, with the gap widening on
long RNA sequences. These results show that EliteDesign improves both design quality and
runtime compared with methods that optimize target folding probability, and allows us to
find high-probability sequences for much longer RNA structures.
